\documentclass[11pt]{article}
\usepackage[letterpaper,margin=0.75in]{geometry}
\usepackage[T1]{fontenc}
\usepackage[utf8]{inputenc}
\usepackage{lmodern}
\usepackage{microtype}
\usepackage{xcolor}
\usepackage{enumitem}
\usepackage{booktabs}

\definecolor{titleline}{RGB}{131,44,39}
\setlength{\parindent}{0pt}
\setlength{\parskip}{0.55em}
\setlist[itemize]{leftmargin=1.2em,itemsep=0.25em,topsep=0.25em}
\setlist[enumerate]{leftmargin=1.4em,itemsep=0.35em,topsep=0.35em}

\begin{document}

{\LARGE\bfseries Scenario 3 Full Example of Play\par}
{\large Harper's Ferry to Port Republic, Six-Turn Historical Walkthrough\par}
\vspace{0.25em}
\color{titleline}\rule{\textwidth}{1.2pt}\par
\vspace{0.45em}

This example presents one plausible six-turn playthrough of Scenario 3, designed to follow the historical flow of Jackson's 1862 Valley campaign as closely as the game system allows. The die rolls are selected to produce a history-leaning result rather than a balanced competitive game. Where the surviving map photographs are not sharp enough to support a fully audited movement proof, the narrative favors historically plausible road corridors and the scenario's actual starting positions.

\textbf{Historical aim of the example.}

\begin{itemize}
\item Jackson strikes Banks first and drives him back instead of lingering in the lower valley.
\item Fremont and Shields are slowed by poor command and distance, so they fail to trap Jackson early.
\item The Confederates use superior concentration and road movement to pivot south.
\item The final turns echo the historical battles of Cross Keys and Port Republic: Ewell holds Fremont while Jackson crushes Shields.
\end{itemize}

\textbf{Starting situation.}

\begin{itemize}
\item Confederate: Jackson, Winder, Ewell, Ashby, and a wagon begin in hex 4511.
\item Union: Banks, Williams, Brodhead, and Detachment E begin in hex 4502; Saxton is in 4611; Fremont, Blenker, and Bayard are in 2516; Shields, Kimball, and Detachment F are in 3519.
\item Scenario 3 lasts 6 turns. The Union Panic Track begins in box 4.
\item In a 1862 scenario the Confederates have the initiative on a die roll of 1--5.
\end{itemize}

\textbf{Turn 1: Jackson batters Banks}

\textbf{Initiative.} Confederate roll \texttt{4}; the Confederates move first.

\textbf{Confederate command and movement.}

Jackson rolls \texttt{5} for command and receives \texttt{3} CPs, enough for \texttt{100\%} movement allowance and a force march. He moves with Winder, Ewell, and Ashby as a single force. Stacking reduces the force by \texttt{1} MP, then force march restores extra MPs equal to Jackson's initiative. The Confederate wagon trails behind on the road rather than entering combat, so the attack remains supplied without slowing the striking force.

Jackson drives straight up the valley road toward Banks at 4502. Banks attempts reaction and fails. Jackson enters the hex, force-march attrition is checked after movement, and the Confederate roll is low enough that no strength is lost.

\textbf{Union command and movement.}

The Union rolls poorly:

\begin{itemize}
\item Banks rolls \texttt{4} for \texttt{2} CPs, giving his hex \texttt{50\%} movement: infantry/leaders \texttt{4} MPs, cavalry \texttt{5} MPs.
\item Fremont rolls \texttt{1} for \texttt{1} CP. With a command rating of \texttt{4}, that gives \texttt{25\%} movement: infantry/leaders \texttt{2} MPs, cavalry \texttt{3} MPs.
\item Shields rolls \texttt{2} for \texttt{1} CP. With a command rating of \texttt{3}, that gives about \texttt{33\%} movement: infantry/leaders \texttt{3} MPs, cavalry \texttt{3} MPs.
\end{itemize}

Banks can maneuver locally, but Fremont and Shields are far too sluggish to support him.

\textbf{Combat.}

Jackson chooses \textit{Pitched Battle}; Banks chooses \textit{Stand}. Jackson's wagon does not deplete. Jackson commits all of Winder, Ewell, and Ashby. Banks commits most of Williams and Brodhead, but not every point available in the hex. Confederate fire is heavy, Union return fire is modest, and the defenders demoralize after the first round. Banks retreats northward, giving Jackson a minor victory and an immediate historical-style shock to the Union position.

\textbf{Conclusion.}

The Panic Track worsens for the Union. Banks's field force is hurt and disordered, while Jackson still has the operational initiative.

\textbf{Turn 2: Jackson breaks contact and turns south}

\textbf{Initiative.} Confederate roll \texttt{3}; the Confederates move first again.

\textbf{Confederate plan.}

Historically, Jackson did not stay to grind against the Potomac line. Instead he slipped away before the Union converging columns could close. The Confederate player follows that logic here. Jackson does not press a second major attack on Banks. Ashby screens northward while Jackson, Winder, and Ewell begin pulling south down the valley road network.

\textbf{Command and movement.}

Jackson receives a solid but not spectacular command result, enough for near-full movement. The wagon follows behind the main body rather than racing ahead. The Confederate force uses road movement to gain distance while staying concentrated.

\textbf{Union half of the turn.}

Banks regains order only partially and remains more concerned with covering the lower valley than sprinting south. Fremont and Shields both get somewhat better command rolls than on Turn 1, but not enough to coordinate a true double envelopment. Fremont starts pushing from the west side of the valley, while Shields edges northward from the Luray side.

\textbf{Operational result.}

By the end of Turn 2, the Union threat is broad but not synchronized. Jackson has traded ground for time and is no longer pinned against Banks.

\textbf{Turn 3: Ashby delays, Jackson concentrates}

\textbf{Initiative.} Confederate roll \texttt{5}; Confederates still move first.

\textbf{Confederate command.}

Jackson gets enough CPs to keep the army moving. Ashby is used aggressively as a cavalry screen, both for reaction threat and for historical flavor. The Confederate player keeps Winder and Ewell close enough to support one another instead of scattering detachments.

\textbf{Movement.}

\begin{itemize}
\item Ashby takes up a delaying position along the western approach, threatening reaction and forcing Fremont to advance carefully.
\item Jackson, Winder, and the wagon continue south through the central valley.
\item Ewell is positioned where he can either support Jackson or turn west to face Fremont if needed.
\end{itemize}

\textbf{Union response.}

This turn the Union command rolls improve, but geography still matters. Fremont can move, but his western route is slower and more awkward. Shields can move, but he is operating on the eastern side and is still not in position to combine directly with Fremont. Banks is no longer the main striking instrument and increasingly becomes a secondary factor.

\textbf{Combat.}

Only light cavalry skirmishing or no decisive battle occurs this turn. That is intentional. Historically, the campaign's middle phase was about delay, concealment, and concentration more than immediate all-out battle.

\textbf{Turn 4: The Union closes, but too late}

\textbf{Initiative.} Confederate roll \texttt{2}; Confederates again seize the first move.

\textbf{Confederate posture.}

Jackson's army is now pivoting toward the Harrisonburg--Cross Keys--Port Republic area. The Confederate player adopts a classic interior-lines stance:

\begin{itemize}
\item Ewell shades west toward Fremont.
\item Winder remains closer to Jackson and the road to Port Republic.
\item The wagon stays central enough to keep the army supplied for a decisive action next turn.
\end{itemize}

\textbf{Union posture.}

The Union player finally has the rough shape of a trap: Fremont from one side, Shields from the other. But the timing still does not work. Their columns are close enough to threaten, not close enough to combine in the same battle on the same turn. That is exactly the strategic failure Jackson historically exploited.

\textbf{Local action.}

Ashby and the Confederate rear guard continue to slow Fremont. Even where no major combat occurs, the threat of reaction, enemy occupation, and terrain costs forces the Union player to spend movement inefficiently.

\textbf{Turn 5: Cross Keys}

\textbf{Initiative.} Confederate roll \texttt{4}; Confederates move first.

\textbf{Confederate design.}

This is the turn where the historical analogy becomes strongest. Ewell is used to face Fremont while Jackson and Winder prepare to strike Shields on the following day. The Confederate player is willing to accept a holding battle in the west if it prevents Union coordination.

\textbf{Command and movement.}

Jackson's roll is adequate, not brilliant. That is enough. Ewell occupies a strong blocking position against Fremont, with Ashby supporting where cavalry can help with reaction and screening.

\textbf{Union command.}

Fremont at last has enough movement to come into contact in strength. Shields also advances, but not in time to combine with Fremont's battle. The Union player must fight piecemeal.

\textbf{Combat at Cross Keys style position.}

Ewell accepts battle against Fremont. The Confederate player chooses a more economical combat option than Jackson used against Banks, because the point now is to hold, not to annihilate. Fremont attacks, but terrain, commitment limits, and Confederate concentration prevent a breakthrough. The Union suffers enough losses and delay that it cannot exploit west of the valley this turn.

\textbf{Operational result.}

By the end of Turn 5, Fremont has been checked. Jackson has bought exactly what he wanted: one more day to fall on Shields without effective interference from the west.

\textbf{Turn 6: Port Republic}

\textbf{Initiative.} Confederate roll \texttt{1}; Confederates still move first at the decisive moment.

\textbf{Confederate plan.}

Now Jackson turns his full striking power against Shields. Ewell is left to contain Fremont as cheaply as possible, while Jackson and Winder mass for the final blow. The wagon is kept close enough to support an all-out attack.

\textbf{Union position.}

Shields is advanced enough to fight, but not supported enough to survive a concentrated Confederate assault. His command roll is serviceable, yet not high enough to produce the kind of movement and commitment needed to escape or to be reinforced in time.

\textbf{Movement and battle.}

Jackson force marches if necessary to close the distance, accepting the attrition risk because this is the last turn and the historical battle demands decisive action. The Confederates strike Shields in a supplied major battle. Jackson chooses \textit{Pitched Battle}; the Union must either stand on unfavorable terms or withdraw and concede the field. In this example Shields stands.

Jackson and Winder commit heavily. Shields commits a respectable portion of his force, but not enough to match Confederate concentration. Confederate fire wins the exchange. The Union either demoralizes during the battle or is forced into a damaging retreat after a sharp final round. The result is a clear Confederate success at the scenario's climax.

\textbf{End of game result.}

The six-turn arc now resembles the historical campaign:

\begin{itemize}
\item Banks is hit first and driven back.
\item Jackson escapes before Fremont and Shields can combine.
\item Fremont is checked in a holding action resembling Cross Keys.
\item Shields is beaten in the final decisive fight resembling Port Republic.
\end{itemize}

That sequence should produce a favorable Panic Track and victory-point outcome for the Confederates, likely a substantive or overwhelming Confederate victory depending on the exact Panic adjustments and battle-victory awards used during play.

\textbf{What this longer example is meant to teach.}

\begin{itemize}
\item Scenario 3 is not just a straight race to battle; it is a campaign of timing, roads, and concentration.
\item Superior Confederate command on the critical turns lets Jackson fight one Union column at a time.
\item Weak Union command does not merely reduce movement. It destroys coordination.
\item The wagon rules matter operationally: keeping supply close without tying the wagon to the lead combat hex is often the right play.
\item The historical Confederate result comes from choosing when \textit{not} to fight just as much as from choosing when to force a major battle.
\end{itemize}

\textbf{Design note.}

This walkthrough is intentionally more narrative than the separate single-turn teaching example. That shorter document is better for learning the exact order of procedures. This longer version is better for understanding how those procedures combine into a historically plausible six-turn campaign.

\end{document}
